% Section 6.2, from lectures/L06/README.md: the goals, the evaluation questions, and the next
% lecture.

\needspace{12\baselineskip}
\section{Review}
\label{c6:sec:review}

\needspace{6\baselineskip}
\subsection*{What you should be able to do now}
After this chapter, you should be able to:
\begin{itemize}
  \item Explain what baud rate is and compute \code{UBRR0} for a given clock and baud rate.
  \item Write a byte, and then a whole string, over UART.
  \item Combine \code{sprintf()} with UART transmission to report a mix of text and numeric data.
  \item Understand why the code waits on \code{UDRE0} before writing a new byte.
  \item Receive a byte over UART and explain how \code{RXC0} mirrors \code{UDRE0}.
\end{itemize}

\needspace{6\baselineskip}
\subsection*{Questions to test yourself}
\begin{enumerate}
  \item What does the \code{UDRE0} flag represent, and what would happen if
    \code{serial_write_byte()} didn't wait for it?
  \item Compute \code{UBRR0} for 19200 baud at a 16\,MHz clock.
  \item Why is \code{sprintf()} used to build a message before transmitting it, rather than trying
    to transmit an \code{int} directly?
  \item What does \code{RXC0} represent, and what would \code{serial_read_byte()} do if it didn't
    wait for it?
\end{enumerate}

\needspace{6\baselineskip}
\subsection*{Looking ahead}
This chapter closes the introductory arc:
\begin{itemize}
  \item C fundamentals (Chapter~\ref{c1:ch}).
  \item Digital I/O and bit manipulation (Chapter~\ref{c2:ch}).
  \item Reacting to events instead of polling for them (Chapter~\ref{c3:ch}).
  \item Reading the analog world (Chapter~\ref{c4:ch}).
  \item Timing without blocking (Chapter~\ref{c5:ch}).
  \item Reporting back to a PC (this chapter).
\end{itemize}
Chapter~\ref{c7:ch} rounds out the peripheral survey with EEPROM, and Chapter~\ref{c8:ch} covers
the watchdog timer. A follow-on arc on driver design starts after that: bundling a peripheral's
state into a struct instead of one-off macros and globals, building directly on the LEDs, buttons,
and timers from these six chapters.
